\chapter{Komposition, Kürzen und Faktorisierung}
\label{b07guide:ch:composition}

\section{Kompositionen beurteilen}

\begin{theorem}[Genaues Surjektivitätskriterium einer Komposition]
\label{b07guide:thm:composition}
Für $f\colon A\to B$ und $g\colon B\to C$ gilt
\[
 \begin{aligned}
 g\circ f\text{ surjektiv auf }C
 &\Longleftrightarrow g[f[A]]=C\\
 &\Longleftrightarrow
   \forall c\in C\quad f[A]\cap g^{-1}[\{c\}]\ne\varnothing.
 \end{aligned}
\]
Gleichwertig ist die Surjektivität von
$g|_{f[A]}\colon f[A]\to C$.
\end{theorem}

\begin{theorem}[Erhalt und Rückschlüsse auf Faktoren]
\label{b07guide:thm:factors}
Seien $f\colon A\to B$ und $g\colon B\to C$.
\begin{enumerate}
\item Sind $f$ und $g$ surjektiv, so ist $g\circ f$ surjektiv.
\item Ist $g\circ f$ surjektiv, so ist $g$ surjektiv.
\item Ist $f$ surjektiv, so gilt
\[
 g\circ f\text{ surjektiv auf }C
 \quad\Longleftrightarrow\quad
 g\text{ surjektiv auf }C.
\]
\end{enumerate}
Aus der Surjektivität der Komposition folgt im allgemeinen
keine Surjektivität des inneren Faktors $f$.
\end{theorem}

\section{Surjektionen rechts kürzen}

\begin{theorem}[Rechtskürzbarkeit als Charakterisierung]
\label{b07guide:thm:cancellation}
Für $f\colon A\to B$ sind äquivalent:
\begin{enumerate}
\item $f$ ist surjektiv auf $B$.
\item Für jede Menge $C$ und alle $u,v\colon B\to C$ gilt
\[
 u\circ f=v\circ f\quad\Longrightarrow\quad u=v.
\]
\item Für eine fest gewählte Menge $E=\{e_0,e_1\}$ mit
      $e_0\ne e_1$ und alle $u,v\colon B\to E$
      gilt dieselbe Kürzungsregel.
\end{enumerate}
Anstelle von $E$ darf jede fest gewählte Menge mit mindestens
zwei verschiedenen Elementen verwendet werden.
Eine Testzielmenge mit höchstens einem Element genügt nicht
als allgemeines Kriterium.
\end{theorem}

\section{Eindeutiger Abstieg auf die Zielmenge}

\begin{definition}[Konstanz auf den Fasern]
Seien $f\colon A\to B$ und $h\colon A\to C$.
Die Funktion $h$ heißt \emph{konstant auf den Fasern von $f$}, wenn
\[
 \forall x,y\in A\quad
 f(x)=f(y)\ \Longrightarrow\ h(x)=h(y).
\]
Eine Funktion $u\colon B\to C$ mit $u\circ f=h$ heißt hier
eine \emph{Faktorisierung von $h$ über $f$}.
\end{definition}

\begin{theorem}[Eindeutige Faktorisierung durch eine Surjektion]
\label{b07guide:thm:factorization}
Seien $f\colon A\twoheadrightarrow B$ und $h\colon A\to C$.
Es gibt genau dann eine Funktion $u\colon B\to C$ mit
\[
 u\circ f=h,
\]
wenn $h$ auf den Fasern von $f$ konstant ist.
Diese Funktion ist dann eindeutig und hat den Graphen
\[
 u=\{(b,c)\in B\times C\mid
           \exists a\in A\ (f(a)=b\land h(a)=c)\}.
\]
Außerdem gilt
\[
 u[B]=h[A],\qquad
 u\text{ surjektiv auf }C
 \ \Longleftrightarrow\
 h\text{ surjektiv auf }C.
\]
Die Aussage gilt auch für leere Mengen.
\end{theorem}

